\begin{table}[htbp]
\centering
\caption{Phillips curve and Okun's law: baseline specifications}
\label{tab:phillips_okun_main}
\begin{threeparttable}
\begin{tabular}{lcc}
\toprule
Specification & $\hat{\beta}$ & $p$-value \\
\midrule
\multicolumn{3}{l}{\textit{Phillips curve} (prices on unemployment)} \\[2pt]
Full panel, contemporaneous & -0.0012*** & $<0.001$ \\
Full panel, city + scenario FE & -0.0012*** & $<0.001$ \\
Baseline scenario, full horizon, contemporaneous & -0.0012*** & $<0.001$ \\
Baseline scenario, full horizon, lagged & -0.0014*** & $<0.001$ \\[4pt]
\multicolumn{3}{l}{\textit{Okun's law} (unemployment on output)} \\[2pt]
Full panel & -0.203** & $0.0103$ \\
Full panel, city + scenario FE & -0.203** & $0.0103$ \\
Baseline scenario, full horizon & -0.191*** & $0.0089$ \\
\bottomrule
\end{tabular}
\begin{tablenotes}
\footnotesize
\item Notes: panel regressions on monthly, city-level simulated series with metropolitan-area fixed effects and standard errors clustered by metropolitan area. The Phillips specification regresses monthly inflation on the unemployment rate; the Okun specification regresses monthly GDP growth on the month-on-month change in unemployment, first-differenced within each simulation so no difference is taken across runs. ``Full panel'' pools all six policy configurations, ``baseline scenario'' restricts to the reference configuration (medium interest rate, improvement policy inactive). ``City + scenario FE'' adds fixed effects for the interest-rate $\times$ improvement-policy cell; ``lagged'' uses $X_{t-1}$ in place of $X_{t}$.
\item The first 18 months of every run are dropped as burn-in: unemployment falls steeply from its initialisation value over that transient, and including it dominates the estimated relationships.
\item Both relationships are negative and significant in every specification, consistent with the target stylised facts described in Section~\ref{sec:calibration}.
\item *** $p<0.01$, ** $p<0.05$, * $p<0.10$.
\end{tablenotes}
\end{threeparttable}
\end{table}